% part: intuitionistic-logic
% chapter: propositions-as-types
% section: type-preservation

\documentclass[../../../include/open-logic-section]{subfiles}

\begin{document}

\olfileid{int}{pty}{tpr}

\olsection{د ټايپ ساتنه}

هغه راکموونکي او ترتيبي بدلونونه چې ومو کتل، د
\Log{tN2i} او \Log{tN3i} نظامونو دپاره هم نېغ په نېغه
تعريفېدلے شي۔ د \Log{tN3i} نظام د سياق~$\Gamma$ په
نسبت د $\lambd$-ترم~$N$ ټايپ هغه !!{formula}~$!A$
تعريفوي چې $\Gamma \Sequent N : !A$ په \Log{tN3i}
کښې !!a{proof} ولري۔ دا چې د $\lambd$-ترمونو د
$\beta$-راکمول قواعد د !!{proof}s راکموونکي او ترتيبي
بدلونونه کټ مټ منعکسوي، يوه مهمه پايله لري: د سياق~$\Gamma$
په نسبت د $\lambd$-ترم ټايپ له $\beta$-بدلون وروسته
هم هماغه پاتې کېږي۔ دې خاصيت ته \emph{د ټايپ ساتنه} وايي۔

\begin{prop}
که $N$ د $\Gamma$ په نسبت ټايپ~$!A$ ولري او
$N \redone N'$ وي، نو $N'$ هم د~$\Gamma$ په نسبت
ټايپ~$!A$ لري۔
\end{prop}

\begin{proof}
  پر~$N$ او د $\Gamma \Sequent N : !A$ د !!{proof}s~$\delta$
  پر لوړوالي په استقرا دا خبره په اسانه ثابتولے شو: که $M$
  د~$N$ يو فرعي ترم وي، نو $\delta$ يو فرعي !!{proof}~$\delta_1$
  لري چې په $\Gamma' \Sequent M : !B$ پاې ته رسېږي۔
  که $M$ رېډکس وي، نو پر~$\delta_1$ يو بدلون پلي کېږي۔
  پر $\delta_1$ د راکمول په پلي کولو د $\Gamma' \Sequent M' : !B$
  !!a{proof}~$\delta_1'$ ترلاسه کوو۔ د \Log{tN3i} د بدلونونو
  له کتنې څرګندېږي چې $M'$ د~$M$ راکم شوے ترم دے۔
  په~$\delta$ کښې د $\delta_1$ پر ځای د $\delta_1'$ په ايښودلو
  د $\Gamma \Sequent N' : !A$ !!a{proof}~$\delta'$ جوړېږي؛
  دلته $N'$ په~$N$ کښې د فرعي ترم~$M$ پر ځای د~$M'$
  په ايښودلو ترلاسه کېږي۔

  مثلاً د~\Intro{\land} هغه راکموونکے بدلون وګورئ
  چې ورپسې \Elim{\land} راځي:
\[
  \bottomAlignProof
  \AxiomC{}
  \Deduce$\Gamma \fCenter N_1 : !B_1$
  \AxiomC{}
  \Deduce$\Gamma \fCenter N_2 : !B_2$
  \RightLabel{\Intro{\land}}
  \BinaryInf$\Gamma \fCenter \pair{N_1}{N_2} : !B_1 \land !B_2$
  \RightLabel{\Elim{\land}[i]}
  \UnaryInf$\Gamma \fCenter \proj{i}{\pair{N_1}{N_2}} : !B_i$
  \DisplayProof
  \quad\redone\quad
  \bottomAlignProof
  \AxiomC{}
  \Deduce$\Gamma \fCenter N_i : !B_i$
  \DisplayProof
  \]
  وينو چې د ښي‌لاس !!{proof} ثبوت ترم، يعنې
  $\delta_1'$، $N_i$ دے۔ دا کټ مټ هغه پايله ده چې د
  کيڼ‌لاس !!{proof}~($\delta_1$) پر ثبوت ترم
  $\proj{i}{\pair{N_1}{N_2}}$ د $\beta$-بدلون له
  پلي کولو ترلاسه کېږي۔

  يا د \Intro{\lif} هغه استنتاج وګورئ چې ورپسې
  د \Elim{\lif} استنتاج راځي، او دا پرې پلي کېدونکے
  راکموونکے بدلون:
\begin{multline*}
  \bottomAlignProof
  \AxiomC{$\delta_2$}
  \Deduce$x: !B_1, \Gamma \fCenter N : !B_2$
  \RightLabel{\Intro{\lif}}
  \UnaryInf$\Gamma \fCenter \lambd[\typeof{x}{!B_1}][N] : !B_1 \lif !B_2$
  \AxiomC{}
  \RightLabel{$\delta_3$}
  \Deduce$\Gamma \fCenter M : !B_1$
  \RightLabel{\Elim{\lif}}
  \BinaryInf$\Gamma \fCenter (\lambd[\typeof{x}{!B_1}][N] M) : !B_2$
  \DisplayProof
  \quad\redone\\
  \bottomAlignProof
  \AxiomC{}
  \RightLabel{$\delta_3$}
  \Deduce$\Gamma \fCenter M : !B_1$
  \RightLabel{$\Subst{\delta_2}{\delta_3}{x:!B_1}$}
  \Deduce$\Gamma \fCenter \Subst{N}{M}{x} : !B_2$
  \DisplayProof
  \end{multline*}
  يو ځل بيا د ښي‌لاس !!{proof} ثبوت ترم کټ مټ د
  کيڼ‌لاس !!{proof} د ثبوت ترم د $\beta$-راکمول پايله ده۔
  البته، دا خبره به په احتياط سره، د~$\delta_2$ پر لوړوالي
  په استقرا، کتل پکار وي: که په~$\delta_2$ کښې د
  $\Gamma' \Sequent x:!B_1$ په بڼه د ټولو اکسيومونو
  پر ځای د $\Gamma \Sequent M:!B_1$ !!{proof}~$\delta_3$
  واچوو، نو په رښتيا د $\Gamma \Sequent \Subst{N}{M}{x} : !B_2$
  !!a{proof} $\Subst{\delta_2}{\delta_3}{x:!B_1}$ جوړېږي۔
\end{proof}

د !!{proof}s او $\lambd$-ترمونو ترمنځ ټينګه مطابقت
بله پايله هم لري۔ که $N$ د سياق~$\Gamma$ په نسبت
د ټايپ $!A$ يو $\lambd$-ترم وي او رېډکس~$M$ ولري
ـ يعنې عادي بڼه نۀ وي ـ نو د $\Gamma \Sequent N : !A$
اړوند \Log{tN3i} ثبوت~$\delta$ يو فرعي !!{proof}
لري چې په $\Gamma' \Sequent M: !B$ پاې ته رسېږي
او بدلون پرې پلي کېدلے شي۔ که $N \redone N'$ د
$M$ پر ځای د هغه د راکم شوي ترم په ايښودلو وي، نو
پر~$\delta$ د راکمول د پلي کولو اړونده پايله د
$\Gamma \Sequent N' : !A$ !!a{proof} ده۔ نو هر
$\lambd$-ترم چې په عادي بڼه کښې نۀ وي، په $\beta$
يو بل ترم ته راکمېږي۔ دې خاصيت ته \emph{پرمختګ} وايي۔

پرمختګ او د ټايپ ساتنه په ګډه دا خبره باوري کوي چې
د $\lambd$-ترمونو $\beta$-راکمول هېڅکله ``نۀ بندېږي'':
که يو ترم صحيح ټايپ ولري او په عادي بڼه کښې نۀ وي، نو
بل صحيح ټايپ لرونکي ترم ته راکمېږي، او د هغه ترم ټايپ
هم هماغه وي۔

\end{document}
